\section{Scope and non-claims}\label{sec:scope}

The results of this paper differ sharply in strength, and their value depends on keeping them
apart. This section collects the boundaries.

\paragraph{Proved here (\inhouse).} The finite tier: the packaging quotient and its
identification with $\HF$ (Theorem~\ref{thm:p5}); the closure theorem and the satisfaction of
every $\mathsf{ZFC}$ axiom except Infinity by $(\HF,\in)$ (Theorems~\ref{thm:hfclosure},
\ref{thm:hfsat}); the rank-audit equivalence (Theorem~\ref{thm:tfae}); non-attainment of
inductive sets by finitely generated rules (Theorem~\ref{thm:nonattain}); the rank-erasure
translation (Theorem~\ref{thm:translation}); the finite Cohen bridge, its density defects and the
counting lemma (Lemma~\ref{lem:meeting}, Theorems~\ref{thm:densitydefect} and \ref{thm:counting});
the infinite-index density facts (Corollary~\ref{cor:infinitecontrast}); the fact that a map on
$\omega$ under which all singletons are visible is injective (Proposition~\ref{prop:nolens}, whose
consequence for $\mathcal B_M$ uses \textbf{J1}); the algebraic strictness
criterion (Theorem~\ref{thm:strict}(1)); and the finite sections together with the standard
$\mathsf{ZF}$ equivalences used to price Choice (Theorem~\ref{thm:choice}). Global statements about
$\HF$ are made in the ambient set-theoretic metatheory. The Lean development verifies selected
cores of these results, listed in \S\ref{sec:mech}.

\paragraph{Proved under named hypotheses (\condg).} The shadow theorem
(Theorem~\ref{thm:shadow}) assumes the structural premise (S) and the commitments
\textbf{C-INF}, \textbf{C-POW}, \textbf{C-REPL} and, for Choice, \textbf{C-SEL}. The forcing
results (Theorems~\ref{thm:novelty}, \ref{thm:recognition}, \ref{thm:strict}(2),
\ref{thm:nonfactor}, Propositions~\ref{prop:converse} and \ref{prop:category}) assume \textbf{J1};
\textbf{J2} is used only to identify the extension $M[G]$. The non-redundancy witnesses
(Proposition~\ref{prop:nonredundancy}(ii)--(iv)) are classical and are used in ambient
$\mathsf{ZFC}$ or in relative-consistency form. Removing these hypotheses removes precisely the
conditional results and leaves the free part standing.

\paragraph{Statements about specific models (\modelg).} The continuum-hypothesis statements
(Theorem~\ref{thm:ch}) are cited classical facts; part (b) needs its own hypothesis of a countable
transitive model of $\mathsf{ZFC}+\mathsf{GCH}$.

\paragraph{What we do \emph{not} claim.}
\begin{enumerate}\itemsep1pt
  \item We do not derive $\mathsf{ZF}$ or $\mathsf{ZFC}$ from finite principles. The infinitary
  content enters through the structural premise and the four commitments, which are hypotheses.
  \item We prove no consistency statement beyond the classical results cited. Nothing here
  implies the existence of a countable transitive model, of an inaccessible cardinal, or the
  consistency of $\mathsf{ZFC}$.
  \item We claim no new independence result. The forcing section reorganizes classical facts and
  adds finite and structural statements about them.
  \item We do not claim that Cohen genericity is the same as non-definability from the ground
  model. Genericity implies novelty and splitting; the converse is false
  (Proposition~\ref{prop:converse}), and genericity is characterized only by the full family of
  coded closed nowhere dense tests (Theorem~\ref{thm:recognition}).
  \item We do not claim that the finite probabilities measure Cohen genericity. They describe
  novelty; Cohen reals have measure zero (Proposition~\ref{prop:category}).
  \item We do not claim that $\mathsf{ZFC}$ is the unique or most natural set theory; the
  continuum-hypothesis statements point the other way.
  \item We make no metaphysical claim about a universe of sets. ``Set'' is used as a package and
  ``$\in$'' as a lens; the readings are structural.
  \item The reading of all this as an account of where the axioms come from
  (\S\ref{sec:discussion}) is interpretation (\metag). It comes after the theorems and is never a
  premise for any of them.
\end{enumerate}

\paragraph{On the G\"odel--Cohen comparison.} The recurring reading of incompleteness and
independence as the saturation and extension faces of one discipline
(Remarks~\ref{rem:saturation} and \ref{rem:notunique}, \S\ref{sec:discussion}) is structural
framing, marked \metag. Each half is backed by theorems---saturation by
Theorem~\ref{thm:nonattain}, extension by Theorems~\ref{thm:novelty} and \ref{thm:strict}---but
the unification of the two is a way of seeing, not a further theorem.
